\section{Hydra maps and their affine word calculus}
\label{sec:hydra-maps}

The 2026 manual is the controlling source for the word \emph{Hydra}.  Earlier
papers use related residue-class maps, but not always the same formal data.
We first type the current object, then isolate an internal inconsistency in
the source definition and state the convention used in the rest of the
edition.

\subsection{Global-field data}

Let $\F$ be either $\Q$ or a rational function field $\F_q(T)$, and let $K$
be a finite extension of $\F$.  Write $\OK$ for the corresponding ring of
integers and $d=[K:\F]$.  If $\ell$ is a place of $K$, let $K_\ell$ be the
completion and $|\cdot|_\ell$ a fixed representative absolute value.  For
$A\in\End_\F(K)$ put
\begin{equation}
\label{eq:place-operator-seminorm}
 \|A\|_\ell=\sup_{x\in K,\,x\ne0}\frac{|Ax|_\ell}{|x|_\ell}
 \in[0,\infty].
\end{equation}
This quantity is finite exactly when $A$ is continuous for the $\ell$-adic
topology, in which case $A$ extends uniquely to a bounded operator on
$K_\ell$.  All placewise operator statements below include this finiteness
hypothesis.  A general $\F$-linear endomorphism need not be continuous at a
chosen single place: for example, a field automorphism can interchange two
places over the same place of $\F$.  Continuity is automatic when the map is
already typed on the selected local factor, as for scalar multiplication.

Let $\Lambda\subsetneq\OK$ be a nonzero ideal.  Its quotient is finite; set
\[
 p=|\OK/\Lambda|\geq2.
\]
Choose an enumeration of the cosets, with the zero coset indexed by $0$:
\[
 \OK/\Lambda=\{\Lambda_0,\ldots,\Lambda_{p-1}\},\qquad
 \Lambda_0=\Lambda.
\]
This enumeration is additional data.  It supplies a bijection from residue
classes to the digit alphabet $\mathcal A_p=\{0,\ldots,p-1\}$; there is no
canonical digit order for a general ideal quotient.

\begin{definition}[affine $\F$-automorphism]
The affine group used by Siegel is the group of affine $\F$-linear
automorphisms of the affine space underlying $K$:
\[
 \Aff_\F(K)=\operatorname{GL}_\F(K)\ltimes K
            =\End_\F(K)^\times\ltimes K.
\]
An element $(r,c)$ acts by $x\mapsto rx+c$, where
$r\in\End_\F(K)^\times$ and $c\in K$.  Its multiplication law is
\begin{equation}
\label{eq:affine-product}
 (r,c)(s,d)=(r\circ s,\;r(d)+c),
\end{equation}
and its inverse is $(r,c)^{-1}=(r^{-1},-r^{-1}c)$.
\end{definition}

The use of $\End_\F(K)$ rather than scalar multiplication by elements of $K$
is essential.  When $[K:\F]>1$, an $\F$-linear map need not be $K$-linear.
The controlling source calls these maps ``affine lattice morphisms.''  That
name is not used as a type assertion here: an arbitrary element of
$\Aff_\F(K)$ need not preserve $\OK$, $\Lambda$, or any lattice in $K$.
Lattice preservation is always stated separately when it is needed.

\begin{definition}[pre-Hydra]
A \emph{pre-Hydra} is a tuple
\[
 \mathcal H=(K,\Lambda,(H_j)_{j\in\mathcal A_p}),
 \qquad H_j=(r_j,c_j)\in\Aff_\F(K).
\]
The $H_j$ are its branches.  The associated piecewise expression, whenever
it is integer-valued on the selected residue classes, is
\begin{equation}
\label{eq:piecewise-hydra}
 H(x)=\sum_{j=0}^{p-1}\ind{x\in\Lambda_j}\bigl(r_jx+c_j\bigr),
 \qquad x\in\OK.
\end{equation}
\end{definition}

\begin{definition}[admissible and exact-integral]
A pre-Hydra is \emph{admissible} when
\begin{equation}
\label{eq:branch-admissibility}
 H_j(\Lambda_j)\subseteq\OK
 \qquad(0\leq j<p).
\end{equation}
It is \emph{exact-integral} when, more strongly,
\begin{equation}
\label{eq:exact-integrality}
 H_j(x)\in\OK\quad\Longleftrightarrow\quad x\in\Lambda_j
 \qquad(x\in\OK,\ 0\leq j<p).
\end{equation}
This is the domain used in Definition 4.5 of the controlling source.  No
field-wide predicate over $x\in K$ is substituted for it.  An admissible
pre-Hydra defines a map $H:\OK\to\OK$ by
\cref{eq:piecewise-hydra}; this is the Hydra map used below.
\end{definition}

\begin{sourceissue}[the two integrality clauses in the manual]
In Definition 4.3 of arXiv:2601.17030v3, the displayed Hydra is required to
satisfy $r_jz+c_j\in\OK$ for every $z\in\OK$ and every $j$.  Definition 4.5
then calls a Hydra integral when $H_j(z)\in\OK$ \emph{if and only if}
$z\in\Lambda_j$, with $z\in\OK$.  For $p\geq2$ these conditions cannot both
hold: after fixing $j$, choose $z\in\OK\setminus\Lambda_j$.  The first clause
puts $H_j(z)$ in $\OK$, while the second forbids it.  The first clause also
fails for the manual's own principal example $T_3$, whose half-integral
branch formulas are integer-valued only on their assigned parity classes.
The minimal condition needed to define $H:\OK\to\OK$ is
\cref{eq:branch-admissibility}; the iff condition
\cref{eq:exact-integrality} is the stronger branch-correctness property used
in the Correspondence Principle.  These two source clauses and their common
integral domain are preserved rather than replaced by a field-wide
predicate.
\end{sourceissue}

\begin{definition}[proper and centered]
A Hydra is
\begin{enumerate}[label=\textup{(\roman*)}]
\item \emph{proper} if $I-r_0\in\End_\F(K)$ is invertible;
\item \emph{centered} if $H_0(0)=0$, equivalently $c_0=0$.
\end{enumerate}
\end{definition}

Properness gives a unique fixed point for the affine branch labelled $0$,
namely $(I-r_0)^{-1}c_0$, but neither makes that point zero nor makes a
translation by it preserve the branch labelled $0$.  The residue-class
permutation must also be tracked.

\begin{proposition}[translation conjugacy and the zero-class branch]
\label{prop:translation-conjugacy-cosets}
Let $a\in\OK$, let $\tau_a(x)=x+a$, and define
$H^{(a)}=\tau_a^{-1}\circ H\circ\tau_a:\OK\to\OK$.  Translation induces the
permutation $\pi_a$ of $\mathcal A_p$ determined uniquely by
\begin{equation}
\label{eq:translation-coset-permutation}
 \Lambda_{\pi_a(j)}=\Lambda_j+a.
\end{equation}
For $x\in\Lambda_j$, the branch of $H^{(a)}$ on $\Lambda_j$ is
\begin{equation}
\label{eq:translated-Hydra-branch}
 H^{(a)}_j(x)
 =\tau_a^{-1}H_{\pi_a(j)}\tau_a(x)
 =r_{\pi_a(j)}x+
   \bigl(r_{\pi_a(j)}a+c_{\pi_a(j)}-a\bigr).
\end{equation}
Admissibility is preserved by this conjugacy, and exact-integrality is
preserved when $H$ is exact-integral.  In the label-dependent terminology of
the definition, $H^{(a)}$ is proper if and only if
$I-r_{\pi_a(0)}$ is invertible.  The branch of $H^{(a)}$ on the zero class
$\Lambda$ is
centered if and only if
\begin{equation}
\label{eq:translation-centering-condition}
 H_{\pi_a(0)}(a)=a.
\end{equation}
In particular, if $H$ is proper and
$a_0=(I-r_0)^{-1}c_0\in\OK$, then conjugation by $\tau_{a_0}$ is guaranteed
to center the zero-class branch when $a_0\in\Lambda$.  If
$a_0\notin\Lambda$, the equation $H_0(a_0)=a_0$ alone gives no such
conclusion; the exact additional condition is
$H_{\pi_{a_0}(0)}(a_0)=a_0$.
\end{proposition}

\begin{proof}
Addition by $a$ is a bijection of $\OK$ carrying each coset
$\Lambda_j$ onto the coset $\Lambda_j+a$.  Since the displayed enumeration
contains every coset exactly once, \cref{eq:translation-coset-permutation}
defines a permutation, with inverse $\pi_{-a}$.  If $x\in\Lambda_j$, then
$x+a\in\Lambda_{\pi_a(j)}$, so the piecewise definition of $H$ selects
$H_{\pi_a(j)}$.  Applying $\tau_a^{-1}$ and expanding the affine map proves
\cref{eq:translated-Hydra-branch}.  If $x\in\Lambda_j$, admissibility of $H$
puts $H_{\pi_a(j)}(x+a)$ in $\OK$, so subtracting $a\in\OK$ proves
admissibility of $H^{(a)}$.  If $H$ is exact-integral, then, for $x\in\OK$,
\[
 H^{(a)}_j(x)\in\OK
 \Longleftrightarrow H_{\pi_a(j)}(x+a)\in\OK
 \Longleftrightarrow x+a\in\Lambda_{\pi_a(j)}
 \Longleftrightarrow x\in\Lambda_j,
\]
which proves exact-integrality.  At $j=0$, centeredness means precisely
$H^{(a)}_0(0)=0$; by \cref{eq:translated-Hydra-branch}, this is equivalent to
$r_{\pi_a(0)}a+c_{\pi_a(0)}-a=0$, which is
\cref{eq:translation-centering-condition}.  The linear part of that same
zero-class branch is $r_{\pi_a(0)}$, proving the properness criterion.  Finally,
$\pi_{a_0}(0)=0$ exactly when $\Lambda+a_0=\Lambda$, equivalently
$a_0\in\Lambda$; in that case the fixed-point equation for $H_0$ supplies
the required equality.
\end{proof}

\subsection{Exact scalar prior-art boundary}

The scalar integer specialization is not a new class of maps under a new
name.  Matthews's generalized Collatz maps already have exactly the same
fixed-modulus residue-class presentation, while Kohl's residue-class-wise
affine maps give the ambient mapping monoid
\cite[Sec.~2]{MatthewsGeneralized}\cite[Defs.~1.1 and~1.4]{KohlRCWA}.
The extra global-field data, non-scalar endomorphisms, centeredness, and
properness in Siegel's formalism remain additional structure.  In scalar
coordinates, exact-integrality is a predicate of the displayed presentation,
whereas Matthews attaches \emph{relatively prime type} to the map at its least
modulus.  The following comparison proves the exact relation without
identifying those two levels prematurely.

\begin{proposition}[fixed-modulus scalar Hydra--Matthews equivalence]
\label{prop:scalar-hydra-matthews}
Fix $D\in\Z_{\geq2}$, put $K=\F=\Q$, $\OK=\Z$, $\Lambda=D\Z$, and use the standard
classes $\Lambda_i=i+D\Z$ for $0\leq i<D$.  Let
\[
 H_i(x)=\alpha_i x+\beta_i,
 \qquad \alpha_i\in\Q^\times,\quad \beta_i\in\Q.
\]
Let $H:\Z\to\Q$ be the associated piecewise function
\[
 H(x)=H_i(x)\quad\text{for }x\in i+D\Z.
\]
Then the following are equivalent for each $i$.
\begin{enumerate}[label=\textup{(\roman*)}]
\item $H_i(i+D\Z)\subseteq\Z$.
\item There are unique integers
      $m_i\ne0$ and $\rho_i$ such that
      \begin{equation}
      \label{eq:matthews-scalar-branch}
       \rho_i\equiv i m_i\pmod D,
       \qquad
       H_i(x)=\frac{m_i x-\rho_i}{D}
       \quad(x\in\Q).
      \end{equation}
\end{enumerate}
Consequently an admissible scalar Hydra with this fixed standard modulus is
exactly a Matthews generalized Collatz presentation with the same displayed
modulus, and its associated map has type $H:\Z\to\Z$.  Write
$d_{\mathrm{Mat}}(H)$ for Matthews's least affine residue modulus, and put
$g_i=\gcd(D,m_i)$.  The full integral locus of the $i$th formal
branch is
\[
 \bigl\{x\in\Z:H_i(x)\in\Z\bigr\}
 =i+\frac{D}{g_i}\Z.
\]
Therefore
\begin{equation}
\label{eq:scalar-exact-integral-relatively-prime}
 \begin{aligned}
 &\text{the scalar Hydra is exact-integral}\\
 &\quad\Longleftrightarrow\quad
   \gcd(D,m_i)=1\ \text{for every }i\\
 &\quad\Longleftrightarrow\quad
   \gcd(m_0\cdots m_{D-1},D)=1.
 \end{aligned}
\end{equation}
The last condition is the relatively-prime numerical condition for this
displayed $D$-presentation.  Matthews assigns the name
\emph{relatively prime type} using the least modulus
$d_{\mathrm{Mat}}(H)$ \cite[Sec.~2]{MatthewsGeneralized}; consequently it is
the same map-level condition only when $D=d_{\mathrm{Mat}}(H)$.

The qualifier is necessary.  The shortened $3x+1$ map has least-modulus
presentation $D=2$ with multipliers $(1,3)$, hence is of Matthews relatively
prime type.  Refining the same map to $D=4$ gives multipliers
$(2,6,2,6)$ and remainder parameters $(0,-2,0,-2)$ in
\cref{eq:matthews-scalar-branch}.  Each formal branch is then integral on a
whole parity class, not on only its assigned class modulo $4$; the refined
$D=4$ presentation is not exact-integral.
\end{proposition}

\begin{proof}
Assume (i).  Since both $i$ and $i+D$ belong to the assigned class,
\[
 m_i:=D\alpha_i=H_i(i+D)-H_i(i)\in\Z.
\]
The branch is an affine automorphism of $\Q$, so $m_i\ne0$.  Put
\[
 \rho_i:=m_i i-DH_i(i)=-D\beta_i.
\]
This is an integer and satisfies
$\rho_i\equiv m_i i\pmod D$; expanding the definition gives
\cref{eq:matthews-scalar-branch}.  The coefficients of an affine function
make $m_i=D\alpha_i$ and $\rho_i=-D\beta_i$ unique.

Conversely, if (ii) holds and $x=i+Dt$, then
\[
 H_i(x)=\frac{m_i i-\rho_i}{D}+m_it\in\Z,
\]
  so (i) holds.  Applying the equivalence independently to all $D$ classes
identifies the two fixed-modulus presentations.

For $x\in\Z$, the congruence in
\cref{eq:matthews-scalar-branch} gives
\begin{align*}
 H_i(x)\in\Z
 &\Longleftrightarrow D\mid m_i x-\rho_i\\
 &\Longleftrightarrow D\mid m_i(x-i)\\
 &\Longleftrightarrow \frac{D}{g_i}\mid x-i.
\end{align*}
This proves the displayed integral locus.  It equals the assigned class
$i+D\Z$ exactly when $g_i=1$.  Requiring this for all $i$ is equivalent to
$\gcd(m_0\cdots m_{D-1},D)=1$, which proves
\cref{eq:scalar-exact-integral-relatively-prime} and the source-term
identification.
\end{proof}

\begin{proposition}[typed inclusion in the rcwa mapping monoid]
\label{prop:scalar-hydra-rcwa}
Inherit $D\in\Z_{\geq2}$ and the branch data from
\cref{prop:scalar-hydra-matthews}, assume its equivalent conditions hold for
every $0\leq i<D$, and define the piecewise map
\[
 H:\Z\longrightarrow\Z,
 \qquad H(x)=H_i(x)\quad\text{when }x\equiv i\pmod D.
\]
Then $H$ is a residue-class-wise affine (rcwa) mapping in Kohl's sense.  If
$d_{\mathrm{Mat}}(H)$ is Matthews's least affine residue modulus and
$\operatorname{Mod}(H)$ is Kohl's least rcwa modulus, then
\begin{equation}
\label{eq:rcwa-modulus-divides-hydra-modulus}
 d_{\mathrm{Mat}}(H)=\operatorname{Mod}(H)\mid D.
\end{equation}
If the displayed scalar Hydra is exact-integral, then
\begin{equation}
\label{eq:exact-integral-minimal-modulus}
 d_{\mathrm{Mat}}(H)=\operatorname{Mod}(H)=D.
\end{equation}
More precisely,
\begin{equation}
\label{eq:scalar-exact-integral-minimal-relative-type}
 \text{the displayed $D$-presentation is exact-integral}
 \quad\Longleftrightarrow\quad
 \left\{
 \begin{array}{l}
  D=d_{\mathrm{Mat}}(H),\\
  H\text{ is of Matthews relatively prime type}.
 \end{array}
 \right.
\end{equation}
Finally,
\[
 H\in\operatorname{RCWA}(\Z)
 \quad\Longleftrightarrow\quad
 H:\Z\to\Z\text{ is bijective},
\]
because $\operatorname{RCWA}(\Z)$ denotes Kohl's group of rcwa
\emph{permutations}, not the full rcwa mapping monoid.
\end{proposition}

\begin{proof}
On the class $i+D\Z$, formula
\cref{eq:matthews-scalar-branch} has Kohl's integer-coefficient form
$(a_i x+b_i)/c_i$ with $(a_i,b_i,c_i)=(m_i,-\rho_i,D)$; hence $D$ is an
rcwa presentation modulus.

Matthews and Kohl minimize over the same positive integers in this scalar
setting.  Matthews asks that the restriction to each residue class be affine;
Kohl writes the same rational-affine restriction as $(ax+b)/c$ with integer
coefficients, which is possible after clearing denominators.  Thus
$d_{\mathrm{Mat}}(H)=\operatorname{Mod}(H)$.

For completeness, let $m=\operatorname{Mod}(H)$ and
$g=\gcd(m,D)$.  Fix a class modulo $g$.  Every class modulo $m$ and every
class modulo $D$ lying above it have a nonempty infinite intersection by the
Chinese remainder theorem.  The affine formulas from the two presentations
agree with $H$ on that intersection, hence agree as rational affine functions.
It follows that all formulas above the fixed class modulo $g$ are identical;
therefore $g$ is also an rcwa presentation modulus.  Minimality of $m$ and
$g\leq m$ force $g=m$, proving $m\mid D$.  The last equivalence is precisely
Kohl's definition of the group $\operatorname{RCWA}(\Z)$.

Now assume exact-integrality.  If $m<D$, then $m\mid D$ and the two distinct
$D$-classes $D\Z$ and $m+D\Z$ lie in the single class $m\Z$.  Let $A$ be the
rational-affine formula of a modulus-$m$ presentation on $m\Z$.  It agrees
with $H_0$ on the infinite set $D\Z$ and with $H_m$ on the infinite set
$m+D\Z$; equality on an infinite subset forces equality of rational-affine
functions, so $H_0=A=H_m$.  Exact-integrality says that the integral locus of
$H_0$ is $D\Z$ and that of $H_m$ is $m+D\Z$.  One function cannot have these
two distinct integral loci, a contradiction.  Hence $m=D$, proving
\cref{eq:exact-integral-minimal-modulus}.  Together with
\cref{eq:scalar-exact-integral-relatively-prime}, this proves the forward
implication in
\cref{eq:scalar-exact-integral-minimal-relative-type}: the displayed modulus
is now the least modulus, and its multipliers satisfy Matthews's defining
coprimality condition.  Conversely, if $D=d_{\mathrm{Mat}}(H)$ and $H$ is of
Matthews relatively prime type, its unique least-$D$ multipliers satisfy
$\gcd(m_0\cdots m_{D-1},D)=1$; the earlier equivalence then proves that the
displayed presentation is exact-integral.
\end{proof}

Thus neither ``Hydra'' nor ``rcwa group element'' may be used as an
unqualified synonym for the other.  The shortened Collatz map itself is an
rcwa mapping which is surjective but not injective
\cite[Rem.~1.2]{KohlRCWA}; it therefore lies outside
$\operatorname{RCWA}(\Z)$.  Conversely, rcwa membership alone supplies no
centeredness, properness, exact-integrality, branch-correctness, or
global-field structure.  The scalar equivalence between exact-integrality and
Matthews's relatively-prime condition also requires that the displayed
modulus be the least one, as proved above; it is not a consequence of rcwa
membership alone.

\subsection{The shortened \texorpdfstring{$qx+1$}{qx+1} Hydra}

For odd $q$, \cref{eq:Tq} is a $2$-Hydra over $K=\F=\Q$ with
$\Lambda=2\Z$ and branches
\begin{equation}
\label{eq:Tq-branches}
 H_0(x)=\frac{x}{2},
 \qquad
 H_1(x)=\frac{qx+1}{2}.
\end{equation}
It is centered and proper because $c_0=0$ and $1-r_0=1/2$.  On integral
inputs it has the exact branch-correctness property
\[
 H_0(x)\in\Z\Longleftrightarrow x\in2\Z,
 \qquad
 H_1(x)\in\Z\Longleftrightarrow x\in1+2\Z
 \quad(x\in\Z).
\]
The unshortened Collatz map is not exact-integral in this sense because its
odd branch $3x+1$ is integer-valued on both parity classes.

\subsection{Words and order}

Let $\Sigma_p^*$ be the free monoid of finite words in $\mathcal A_p$,
including the empty word $\varnothing$.  Concatenation is denoted by
$\mathbf i\wedge\mathbf j$.  For
$\mathbf j=(j_0,\ldots,j_{m-1})$, set
\begin{equation}
\label{eq:word-composition}
 H_{\mathbf j}=H_{j_0}\circ H_{j_1}\circ\cdots\circ H_{j_{m-1}},
 \qquad H_\varnothing=\operatorname{id}_K.
\end{equation}
Thus
\begin{equation}
\label{eq:concat-composition}
 H_{\mathbf i\wedge\mathbf j}=H_{\mathbf i}\circ H_{\mathbf j}.
\end{equation}
There are unique maps
\[
 M_H:\Sigma_p^*\to\End_\F(K)^\times,
 \qquad X_H:\Sigma_p^*\to K
\]
such that
\begin{equation}
\label{eq:affine-word-decomposition}
 H_{\mathbf j}(x)=M_H(\mathbf j)x+X_H(\mathbf j).
\end{equation}
The second component is Siegel's \emph{numen}; equivalently,
$X_H(\mathbf j)=H_{\mathbf j}(0)$.

\begin{proposition}[affine concatenation identities]
\label{prop:affine-concatenation}
For all $\mathbf i,\mathbf j\in\Sigma_p^*$,
\begin{align}
 M_H(\mathbf i\wedge\mathbf j)
   &=M_H(\mathbf i)M_H(\mathbf j),\label{eq:M-concat}\\
 X_H(\mathbf i\wedge\mathbf j)
   &=M_H(\mathbf i)X_H(\mathbf j)+X_H(\mathbf i)
     =H_{\mathbf i}\bigl(X_H(\mathbf j)\bigr).
     \label{eq:X-cocycle}
\end{align}
\end{proposition}

\begin{proof}
Substitute \cref{eq:affine-word-decomposition} twice into
\cref{eq:concat-composition}:
\[
 H_{\mathbf i}(H_{\mathbf j}(x))
 =M_H(\mathbf i)M_H(\mathbf j)x
  +M_H(\mathbf i)X_H(\mathbf j)+X_H(\mathbf i).
\]
Uniqueness of the linear and translation parts gives both identities.
\end{proof}

\begin{corollary}[closed word formulas]
\label{cor:closed-word-formulas}
For $\mathbf j=(j_0,\ldots,j_{m-1})$,
\begin{align}
 M_H(\mathbf j)
   &=r_{j_0}r_{j_1}\cdots r_{j_{m-1}},\label{eq:M-word}\\
 X_H(\mathbf j)
   &=\sum_{a=0}^{m-1}
       \bigl(r_{j_0}r_{j_1}\cdots r_{j_{a-1}}\bigr)c_{j_a},
       \label{eq:X-word}
\end{align}
where an empty product is $I$ and products mean composition in the displayed
order.
\end{corollary}

\begin{proof}
Induct on $m$.  Appending the last letter and applying
\cref{eq:X-cocycle} gives the new summand after the existing left prefix;
\cref{eq:M-concat} gives the product formula.
\end{proof}

When the $r_j$ do not commute, the product order in
\cref{eq:M-word,eq:X-word} is indispensable.  Expressions such as
$(1-M_H(\mathbf j))^{-1}$ mean the inverse of an endomorphism of $K$, not
scalar division.

\subsection{Repeated words and affine fixed points}

For a word $\mathbf j$, write $\mathbf j^{\wedge n}$ for its $n$-fold
concatenation, and abbreviate $M=M_H(\mathbf j)$ and $X=X_H(\mathbf j)$.

\begin{proposition}[geometric sum in the endomorphism algebra]
\label{prop:repeated-word}
For $n\geq1$,
\begin{align}
 M_H(\mathbf j^{\wedge n})&=M^n,\\
 X_H(\mathbf j^{\wedge n})&=(I+M+\cdots+M^{n-1})X.\label{eq:operator-geometric}
\end{align}
If $I-M$ is invertible, then
\begin{equation}
\label{eq:operator-geometric-inverse}
 X_H(\mathbf j^{\wedge n})=(I-M^n)(I-M)^{-1}X.
\end{equation}
The affine map $H_{\mathbf j}$ has a unique fixed point precisely when
$I-M$ is invertible, in which case
\begin{equation}
\label{eq:word-fixed-point}
 \operatorname{Fix}(H_{\mathbf j})=\{(I-M)^{-1}X\}.
\end{equation}
\end{proposition}

\begin{proof}
The first two identities follow by induction from
\cref{prop:affine-concatenation}.  Because $M$ commutes with every polynomial
in $M$,
\[
 (I-M)(I+M+\cdots+M^{n-1})=I-M^n,
\]
which gives \cref{eq:operator-geometric-inverse}.  Finally,
$H_{\mathbf j}(z)=z$ is exactly $(I-M)z=X$.
\end{proof}

\begin{remark}
Properness requires $I-r_0$ to be invertible.  It does not imply that
$I-M_H(\mathbf j)$ is invertible for every nonempty word.  The latter is a
separate hypothesis whenever \cref{eq:word-fixed-point} is used.
\end{remark}

\subsection{Correct words versus formal words}

Every word defines an affine map on $K$, but it need not describe a genuine
orbit segment of the piecewise map $H$.  A chronological branch word
$\boldsymbol\varepsilon=(\varepsilon_0,\ldots,\varepsilon_{m-1})$ is
\emph{correct at $z\in\OK$} if
\[
 H^{\circ k}(z)\in\Lambda_{\varepsilon_k}
 \qquad(0\leq k<m).
\]
Then
\begin{equation}
\label{eq:chronological-composition}
 H^{\circ m}(z)
 =H_{\varepsilon_{m-1}}\circ\cdots\circ H_{\varepsilon_0}(z)
 =H_{\operatorname{rev}(\boldsymbol\varepsilon)}(z).
\end{equation}
Exact-integrality is designed to recover correctness from integrality of
successive branch values.  The passage from a formal fixed point of
$H_{\mathbf j}$ to a periodic point of $H$ therefore requires more than the
affine equation alone; the needed arithmetic step is isolated in
\cref{sec:correspondence}.
